% Two River Bio working note on the trbdoc class (trb-branding, latex/trbdoc v0.6.0).
% Compile:  latexmk note.tex        (latexmkrc selects LuaLaTeX and finds ./trbdoc)
% Audience: `internal' = machine logo + confidential footer; `external' = horizontal logo + (c) footer.
\documentclass[note,internal]{trbdoc}

% Optional extras; delete what you do not use. Do not load amssymb or amsfonts.
\usepackage{tikz}
\usetikzlibrary{arrows.meta,positioning}
\usepackage{multicol}
\usepackage{siunitx}

% --- Title block ---------------------------------------------------------------
\title{Note title}
\subtitle{One line that says what this note decides, proposes or reports}
\doctype{Working note}            % small label above the title and in the footer
\author{Your Name}
\date{20 October 2026}

% --- Your macros ---------------------------------------------------------------
\newcommand{\E}{\mathbb E}
\newcommand{\R}{\mathbb R}

\begin{document}
\maketitle

% Open with the claim. Two to four sentences: what the question was, what the answer is,
% and what should happen next. A reader who stops here should know the outcome.
State the claim first, then the evidence.\cite{example2026}

\paragraph{Context.}
Why this note exists: the question, the decision it feeds, and what was already known.
Run-in heads carry the argument; write \verb|\paragraph{Heading.}| and continue the
sentence. Use \verb|\section| only when the note grows long enough to need navigation.

\paragraph{What we did.}
Method, data and settings, in enough detail to repeat. Put numbers in a table rather
than in prose.

{\small
\renewcommand{\arraystretch}{1.12}
\begin{tabularx}{\linewidth}{@{}>{\raggedright\arraybackslash}p{3.2cm}X>{\raggedleft\arraybackslash}p{2cm}@{}}
\toprule
\textbf{Item} & \textbf{Description} & \textbf{Value}\\
\midrule
Row one & What it is and where it came from & \num{0.00}\\
Row two & What it is and where it came from & \num{0.00}\\
\bottomrule
\end{tabularx}}

\paragraph{What we learned.}
Results and their interpretation. Say what surprised you and what did not.
Math is set in Concrete Math to match the text:
\[
  \hat\theta = \arg\max_{\theta}\; \E_{(x,y)\sim\mathcal D}\bigl[\log p_\theta(y\mid x)\bigr],
  \qquad x\in\R^k.
\]

\paragraph{Open questions.}
\begin{itemize}
  \item What is still unresolved, and what evidence would settle it.
  \item Use \textcolor{trbAccent}{\texttt{trbAccent}} (Splash) only for the one thing that
    must not be missed; everything else stays in the text color.
\end{itemize}

\paragraph{Next step.}
The single concrete action this note commits to, who owns it, and by when.

% Optional figure. Colors by name only: trbHeading, trbLink, trbAccentBlue, trbMuted,
% trbRule, trbSurface, trbAccent, and every ramp shade (trbHudson500, trbEast700, ...).
\begin{figure}[htbp]
\centering
\begin{tikzpicture}[
  box/.style={draw=trbRule,fill=trbSurface,rounded corners=2pt,align=center,
    minimum height=8mm,text width=2.6cm,inner sep=3pt,font=\small},
  arr/.style={-{Latex[length=2mm]},draw=trbLink,line width=0.8pt}]
\node[box] (a) {Step one};
\node[box,right=5mm of a] (b) {Step two};
\node[box,right=5mm of b] (c) {Step three};
\draw[arr] (a)--(b); \draw[arr] (b)--(c);
\end{tikzpicture}
\caption{Caption that states what the figure shows, not just what it is.}
\end{figure}

% References: two compact columns. Delete this block if the note cites nothing.
\begingroup
\renewcommand{\bibsection}{}
\setlength{\columnsep}{16pt}
\begin{multicols}{2}[\paragraph{References}]
\raggedright
\bibliographystyle{unsrtnat}
\bibliography{references}
\end{multicols}
\endgroup
\end{document}
